\section{The full--partial reconstruction mapping}\label{sec:fprm}

\begin{figure*}
\includegraphics[width=\textwidth]{figures/fig_fprm_scheme.pdf}
\caption{The reconstruction problem on the chaotic attractor ($\mu=-0.8$, $\gamma=0.25$, $A=0.8$, $\Omega=2/3$). (a) Measurable observables over six forcing periods: the complete spring-force channel $s$ and the collar-height channel $c$ with $60\%$ of its entries missing (open circles); lines are the continuous trajectories and dots the noisy samples. (b)--(d) Model-free identifiability test (task labels as in Sec.~\ref{sec:tasks}): target value at time $t$ versus its value at the time $t^{\ast}$ of the nearest neighbor of the delay vector of the complete channel (logarithmic histograms of $4\times10^{4}$ pairs, with temporal neighbors closer than one forcing period excluded). (b) $s$ complete, $c$ target: identifiable. (c) $c$ complete, $s$ target: the reflection symmetry creates a mirror branch. (d) As (c), with the forcing phase appended to the delay vector: identifiability is restored.}
\label{fig:fprm_scheme}
\end{figure*}

The attractor of Sec.~\ref{sec:attractor} is never observed as such: sensors record scalar functions of the state, such as the spring force or the height of the collar, and some of these records are incomplete. The reconstruction method is therefore developed in the order dictated by the three properties identified there. The delay reconstruction, which replaces the unobserved state by the history of one complete channel, is introduced first (Sec.~\ref{sec:takens}), followed by the FPRM, which maps this history onto the missing channel (Sec.~\ref{sec:fprm_framework}). The identifiability conditions imposed by the symmetry are then derived and tested on data (Sec.~\ref{sec:identifiability}), and the generative engine and the physical constraints proposed here are presented (Secs.~\ref{sec:flow} and \ref{sec:physics_reg}), before the reconstruction tasks and the metrics are defined (Sec.~\ref{sec:tasks}). In brief, task C1 reconstructs the collar height and the angular velocity from the complete spring force on the chaotic attractor; task C2 reconstructs the spring force from the complete collar height of an undriven, stochastic PIAR; task C3 does so on the chaotic attractor, with and without the forcing phase; and task C0 is a trivial control.

\subsection{Delay reconstruction}\label{sec:takens}

The state $\mathbf z(t)$ of a dynamical system is assumed to settle on an attractor $\mathcal A$ and to be observed through a scalar measurement $x(t)=\mathcal O(\mathbf z(t))$, such as the spring force. The delay vector
\begin{equation}
\mathbf x(t) = \bigl[x(t),\,x(t-\tau),\,\dots,\,x(t-(d_E-1)\tau)\bigr]\in\mathbb R^{d_E},
\label{eq:delay}
\end{equation}
built from $d_E$ samples separated by a delay $\tau$, defines the delay map $\Phi_x:\mathbf z(t)\mapsto\mathbf x(t)$. Takens' theorem~\cite{Takens1981}, in the form extended to fractal attractors by Sauer, Yorke, and Casdagli~\cite{Sauer1991}, states that the delay map is one-to-one on the attractor for a generic measurement function $\mathcal O$ whenever $d_E>2D_0$, where $D_0$ is the box-counting dimension of $\mathcal A$; the extension to periodically forced systems, in which the phase of the drive is part of the state, is due to Stark~\cite{Stark1999}. The reconstructed set $\mathcal A_x=\Phi_x(\mathcal A)$ is then a faithful copy of the original attractor: distinct states have distinct delay vectors, and the history of a single variable encodes the whole state~\cite{Packard1980,Abarbanel1993}. In practice the delay is chosen at the first minimum of the average mutual information between $x(t)$ and $x(t-\tau)$~\cite{FraserSwinney1986}, and the dimension by the method of false nearest neighbors~\cite{Kennel1992}, which increases $d_E$ until points that are neighbors in $\mathbb R^{d_E}$ remain neighbors in $\mathbb R^{d_E+1}$~\cite{KantzSchreiber2004}.

\subsection{The mapping between full and partial reconstructions}\label{sec:fprm_framework}

The FPRM~\cite{Wu2026FPRM} considers two observables of the same system: a completely observed channel $x$ and a partially observed channel $y$ with missing entries. Their delay reconstructions, $\mathcal A_F=\Phi_x(\mathcal A)$ (full) and $\mathcal A_P=\Phi_y(\mathcal A)$ (partial), are both one-to-one images of the attractor, so that the map
\begin{equation}
\Psi = \Phi_y\circ\Phi_x^{-1}:\ \mathcal A_F\to\mathcal A_P
\label{eq:fprm_map}
\end{equation}
is well defined and continuous. In particular, the current value of the incomplete channel is a function of the delay vector of the complete one,
\begin{equation}
y(t) = \psi\bigl(\mathbf x(t)\bigr),
\label{eq:fprm_function}
\end{equation}
where $\psi$ is the first component of $\Psi$. The function $\psi$ can be learned from the pairs $(\mathbf x(t),y(t))$ at the times at which $y$ is observed and evaluated at the times at which it is missing. This is the essence of the method, which is closely related to the cross mapping between reconstructed manifolds used to detect causal coupling~\cite{Sugihara2012}.

In the original formulation, $\psi$ is learned by Gaussian-process regression~\cite{Rasmussen2006}. A zero-mean Gaussian-process prior with covariance function $\mathcal K(\mathbf x,\mathbf x')$ is placed on $\psi$, and, given $N$ training inputs with noisy outputs $\mathbf y=(y_1,\dots,y_N)^{\top}$, the prediction at a new input $\mathbf x_*$ is the Gaussian distribution with mean and variance
\begin{equation}
\begin{aligned}
\bar\psi(\mathbf x_*) &= \mathbf k_*^{\top}(\mathbf K+\sigma_n^{2}I)^{-1}\mathbf y,\\
\sigma_*^{2}(\mathbf x_*) &= \mathcal K(\mathbf x_*,\mathbf x_*) - \mathbf k_*^{\top}(\mathbf K+\sigma_n^{2}I)^{-1}\mathbf k_* .
\end{aligned}
\label{eq:gpr}
\end{equation}
Here $\mathbf K$, with entries $\mathcal K(\mathbf x_i,\mathbf x_j)$, is the covariance matrix of the training inputs, $\mathbf k_*$ is the vector of covariances between $\mathbf x_*$ and the training inputs, $\sigma_n^{2}$ is the noise variance, and $I$ is the identity matrix; the mean $\bar\psi$ is the imputed value, and $\sigma_*$ its uncertainty. With a squared-exponential covariance and one length scale per delay coordinate, the hyperparameters are fitted by maximizing the marginal likelihood of the training data. In the implementation used here, the fitted noise variance is added to $\sigma_*^{2}$, so that the predictive distribution refers to a new noisy observation; when the marginal likelihood attributes most of the variance to noise, the prediction becomes broad. The approach is elegant and data-efficient, and it recovers complex dynamics at large missing fractions~\cite{Wu2026FPRM}. Its predictive distribution, however, is a single Gaussian, and this restriction becomes decisive whenever Eq.~\eqref{eq:fprm_function} fails.

\subsection{Identifiability and symmetry}\label{sec:identifiability}

Equation~\eqref{eq:fprm_function} rests on the assumption that the delay map of the complete channel is one-to-one, and symmetries can violate it~\cite{LetellierAguirre2002,Letellier2005}. If a transformation $R$ of the state space commutes with the flow, maps the attractor onto itself, and leaves the complete channel invariant, $x\circ R=x$, but changes the incomplete one, $y\circ R\ne y$, then a state and its image under $R$ produce the same delay vector while carrying different values of $y$. The missing value is then not a function of $\mathbf x(t)$, and its conditional distribution $p(y\mid\mathbf x)$ has one branch for each member of the symmetry orbit. The PIAR realizes this situation in its simplest form. By Eq.~\eqref{eq:reflection}, the collar height $c$ is invariant under the reflection $(\theta,\dot\theta)\to(-\theta,-\dot\theta)$ while the spring force $s$ changes sign, so that in the autonomous system no record of $c$ can determine the sign of $s$, and the conditional distribution of $s$ is bimodal, with branches $s=\pm(1-c^{2})^{1/2}$. The converse task is well posed: $s$ is invariant only under $\theta\to\pi-\theta$, which exchanges the upper and lower half-planes and is not a symmetry of the dynamics, because gravity distinguishes up from down. In the driven system the reflection is a symmetry only when combined with a half-period time shift, Eq.~\eqref{eq:st_symmetry}. Two states related by this symmetry experience forcing phases that differ by $\pi$, and appending the phase $(\cos\Omega t,\sin\Omega t)$ to the delay vector therefore restores the one-to-one property. Stark's forced embedding theorem~\cite{Stark1999} does not require the phase to be observed, but it holds only for generic observables, and $c$ is not generic because it is invariant under Eq.~\eqref{eq:st_symmetry}.

These statements can be tested directly on data, before any model is trained, with the simplest possible estimate of $\psi$: the value of the incomplete channel at the nearest neighbor of $\mathbf x(t)$ on $\mathcal A_F$. Tests of this kind, which ask whether neighbors in one reconstruction remain neighbors in another, have long been used to detect continuity and generalized synchronization between time series~\cite{Pecora1995,Rulkov1995}. Figure~\ref{fig:fprm_scheme} shows the result for the chaotic attractor of Sec.~\ref{sec:attractor}, computed from the complete channel, with measurement noise of standard deviation $0.01$, the noise-free target channel, and the embeddings selected in Sec.~\ref{sec:tasks}; neighbors closer in time than one forcing period were excluded, as recommended by Theiler~\cite{Theiler1986}. When the spring force is complete and the collar height is the target, the neighbor predicts the target almost perfectly (correlation $\rho=0.99$), and the scatter collapses onto the diagonal [Fig.~\ref{fig:fprm_scheme}(b)]. When the roles are exchanged, the scatter forms an X: the neighbor carries either the correct value or its mirror image with equal probability, and the correlation vanishes [$\rho=-0.01$; Fig.~\ref{fig:fprm_scheme}(c)]. Appending the forcing phase removes the mirror branch and restores $\rho=0.99$ [Fig.~\ref{fig:fprm_scheme}(d)]. The identifiability of every reconstruction task is thus established in advance, from the symmetry of the model and from a model-free test.

\subsection{Generative engine: A conditional normalizing flow}\label{sec:flow}

To represent ambiguous targets, the FPRM is recast here as conditional density estimation: instead of a function $\psi$, a conditional distribution $q_\phi(y\mid\mathbf h)$ with parameters $\phi$ is learned, and imputations are drawn from it. For the sample recorded at time $t_n$, the conditioning vector
\begin{equation}
\begin{split}
\mathbf h_n = \bigl[\,&x_{n-m\tau},\dots,x_{n+m\tau};\ \tilde y_{n\pm1},\dots,\tilde y_{n\pm n_b};\\
&\mathfrak m_{n\pm1},\dots,\mathfrak m_{n\pm n_b};\ \cos\Omega t_n,\ \sin\Omega t_n\,\bigr]
\end{split}
\label{eq:context}
\end{equation}
collects all the information available at that sample, with subscripts counting samples and the delay $\tau$ expressed in samples. Its first block is a symmetric window of the complete channel with $2m+1$ samples spaced by $\tau$, where $m=\lceil(d_E-1)/2\rceil$; it is a delay vector of dimension $2m+1\ge d_E$ centered at $t_n$, which is admissible because the reconstruction is performed offline. The second and third blocks contain the neighbors of the incomplete channel, $n_b=3$ samples on each side, set to zero where missing ($\tilde y$), together with the corresponding availability mask $\mathfrak m$; they allow the model to exploit continuity when some neighbors are observed. The last block, the forcing phase, is included in task C1 and in the phase variant of task C3 (Sec.~\ref{sec:tasks}).

The conditional density is represented by a normalizing flow~\cite{Papamakarios2021}, which writes the target, of dimension $d_y$ (one or two), as an invertible transformation of a standard normal variable, $y=\mathcal G_\phi(\mathbf u;\mathbf h)$ with $\mathbf u\sim\mathcal N(\mathbf 0,I)$ in $\mathbb R^{d_y}$. The change-of-variables formula then gives the exact density,
\begin{equation}
\ln q_\phi(y\mid\mathbf h) = \ln\mathcal N\bigl(\mathcal G_\phi^{-1}(y;\mathbf h)\bigr) + \ln\left|\det\frac{\partial \mathcal G_\phi^{-1}(y;\mathbf h)}{\partial y}\right| ,
\label{eq:flow}
\end{equation}
in which the first term is the base density evaluated at the preimage of $y$ and the second accounts for the local stretching of volumes by the transformation. A neural spline flow (NSF)~\cite{Durkan2019} is used, in which $\mathcal G_\phi$ is a composition of three autoregressive transforms, each a monotonic rational-quadratic spline with eight bins whose knots are produced by a neural network that receives the conditioning vector; the implementation relies on \texttt{zuko}~\cite{zuko} and PyTorch~\cite{Paszke2019}. The parameters are trained by maximizing the log-likelihood of Eq.~\eqref{eq:flow} over the observed pairs $(\mathbf h_n,y_n)$, missing entries are imputed by sampling, and the posterior median is used whenever a point estimate is required.

A flow was chosen as the main engine for three reasons, two of which are quantified against the diffusion model in Sec.~\ref{sec:res_engine}. It provides the exact likelihood of Eq.~\eqref{eq:flow}, a proper criterion for selecting the model on held-out data, whereas the conditional variational autoencoder and the diffusion model introduced below are validated through a bound or a surrogate of their likelihood; it generates each sample with a small, fixed number of network evaluations instead of an iterative denoising chain; and its samples are differentiable functions of the parameters, obtained in the same pass, which makes the physical constraints described next affordable. One limitation should be anticipated: because a flow is a diffeomorphism of $\mathbb R^{d_y}$, a flow with a unimodal Gaussian base must connect separated modes by thin bridges of probability~\cite{Cornish2020}.

For comparison, two other conditional generative models receive the same conditioning vector: a conditional variational autoencoder (cVAE)~\cite{Kingma2014,Sohn2015}, which maximizes a lower bound on the likelihood through a latent Gaussian variable, and a conditional denoising diffusion probabilistic model (DDPM)~\cite{Ho2020,NicholDhariwal2021}, which learns to invert a gradual noising process and underlies the conditional score-based diffusion imputer (CSDI)~\cite{Tashiro2021}.

\subsection{Physics-informed regularization}\label{sec:physics_reg}

The equation of motion, Eq.~\eqref{eq:forced}, is known in form, and the question is whether imposing it on the imputations improves them. The known mechanics enters, as in physics-informed neural networks~\cite{Raissi2019}, through residuals evaluated on reparameterized samples $\hat y\sim q_\phi(\cdot\mid\mathbf h_n)$~\cite{Kingma2014}, so that gradients propagate through the flow. For the task in which $s$ is complete and the collar height $c$ and the angular velocity $v=\dot\theta$ are imputed (task C1 of Sec.~\ref{sec:tasks}), three residuals vanish for the true values:
\begin{align}
r_{\mathrm{hol}} &= \hat c^{2}+s^{2}-1, \label{eq:res_hol}\\
r_{\mathrm{kin}} &= \hat c\,\hat v - \dot s, \label{eq:res_kin}\\
r_{\mathrm{EOM}} &= \ddot s + s\,\hat v^{2} - \hat c\,\bigl[\beta s - \lambda s\hat c - \gamma\hat v + f(t)\bigr], \label{eq:res_eom}
\end{align}
where $\hat c$ and $\hat v$ are the imputed values and $f(t)=A\cos\Omega t$ is the drive. The holonomic residual, Eq.~\eqref{eq:res_hol}, states that $(s,c)$ lies on the unit circle. The kinematic residual, Eq.~\eqref{eq:res_kin}, follows from differentiating $s=\sin\theta$, which gives $\dot s=\cos\theta\,\dot\theta$. The dynamical residual, Eq.~\eqref{eq:res_eom}, is the equation of motion multiplied by $\cos\theta$ and rewritten with the identity $\ddot s=\cos\theta\,\ddot\theta-\sin\theta\,\dot\theta^{2}$, so that it contains only the observed channel, its derivatives, and the imputed quantities. For the driven task in which $c$ is complete and $s$ is imputed (C3), the holonomic residual $\hat s^{2}+c^{2}-1$ is complemented by the equation of motion expressed through $c$, its derivatives, and $\hat s$, multiplied by $\sin^{3}\theta$ to remove the singular factor $1/\sin\theta$ of $\dot\theta=-\dot c/\sin\theta$ (Appendix~\ref{app:derivations}),
\begin{equation}
r_{\mathrm{EOM}} = -\hat s^{2}\ddot c - c\,\dot c^{2} - \gamma\dot c\,\hat s^{2} + \lambda\hat s^{4}c - \beta\hat s^{4} - f(t)\,\hat s^{3}.
\label{eq:res_c3}
\end{equation}
The same residual, which contains the known drive $f(t)$ at the true sample times, is used in both variants of C3, so that in the variant without the phase the physics term carries timing information that the conditioning vector lacks. In the stochastic task C2 only the holonomic residual is imposed, because the random force of the Langevin dynamics has no deterministic counterpart in the equation of motion.

Derivatives of the complete channel are estimated with Savitzky--Golay filters~\cite{SavitzkyGolay1964} (window of 7 samples, degree 4). Residuals computed from noisy derivatives cannot vanish exactly, and forcing them to zero oversharpens the conditional density. Each residual $r_j$, with $j\in\{\mathrm{hol},\mathrm{kin},\mathrm{EOM}\}$, is therefore scaled by its root-mean-square value $\sigma_j$ at the observed data, and only the part with $|r_j|>2\sigma_j$ is penalized, which yields the training objective
\begin{equation}
\begin{split}
J(\phi) = &-\bigl\langle\ln q_\phi(y_n\mid\mathbf h_n)\bigr\rangle_{\mathrm{obs}}\\
&+ w\,\Bigl\langle\sum_j \mathrm{ReLU}\bigl(r_j^{2}/\sigma_j^{2}-4\bigr)\Bigr\rangle_{\mathrm{coll},\,\hat y}.
\end{split}
\label{eq:loss}
\end{equation}
The first average runs over the observed pairs and is the negative log-likelihood; the second runs over collocation times and samples, $\mathrm{ReLU}(\cdot)=\max(\cdot,0)$, and the weight $w=1$ is ramped in linearly between $10\%$ and $20\%$ of the training budget. The collocation times are drawn from all time indices, including those whose target is missing: because the residuals depend only on the complete channel and on the samples, this adds information precisely where imputation is required, without using any held-out value. One property of these residuals is decisive for what follows. In tasks C2 and C3 the holonomic residual is invariant under $\hat s\to-\hat s$, and in C3 only the forcing term of Eq.~\eqref{eq:res_c3} is odd in $\hat s$, so that the physics can discriminate between mirror branches only through the drive. In task C1 the relevant mirror is $(\hat c,\hat v)\to(-\hat c,-\hat v)$, which leaves $s$ and $\dot s$ unchanged; under it, only the gravitational and forcing terms of Eq.~\eqref{eq:res_eom} change sign, so that the physics can locate the rod in the upper or lower half-plane.

\subsection{Reconstruction tasks, data, and metrics}\label{sec:tasks}

The tasks follow from the identifiability analysis of Sec.~\ref{sec:identifiability}. Reconstructing $\dot\theta$ from a complete record of $\theta$ is solved by numerical differentiation and is retained only as a control (C0). Three informative tasks use the measurable observables $s$ and $c$. In task C1 the spring force is complete, and the collar height and the angular velocity are missing at the same instants, as happens when both are derived from a single angle-encoder record; this target is identifiable, and the angular velocity cannot be obtained by differentiating the observed channel alone. In task C2 the collar height of an autonomous, stochastic PIAR is complete and the spring force is missing, so that the target is not identifiable by symmetry. Task C3 is the driven analogue of C2, with and without the forcing phase in the conditioning vector.

Tasks C1 and C3 use the reference attractor of Sec.~\ref{sec:attractor}, sampled 32 times per forcing period ($\Delta t=0.295$) for $4\times10^{4}$ samples after a transient of 400 periods; in C1 the forcing phase is part of the conditioning vector. Task C2 uses a Langevin trajectory integrated with the BAOAB scheme~\cite{LeimkuhlerMatthews2013} for $\mu=-0.5$, $\gamma=0.1$, and temperature $T_b=0.6\,MgL$, sampled every $0.2$ time units; the rod librates about the hanging position and occasionally escapes over the upright barrier into rotation, at a rate governed by Kramers' theory of thermally activated escape~\cite{Kramers1940}. For such a stochastically forced system, delay embeddings remain meaningful in the generalized sense of Stark et al.~\cite{Stark2003}, but the false-nearest-neighbor fraction never falls below $7\%$, and the dimension with the smallest fraction was used; the same rule applies to C3, whose smallest fraction is $1.3\%$. Independent Gaussian noise of standard deviation $0.01$ is added to every channel, and entries of the incomplete channel are removed completely at random~\cite{LittleRubin2019} with fractions of $30\%$, $60\%$, and $90\%$, with two random masks per fraction. The delays and dimensions selected by mutual information and false nearest neighbors are $(\tau,d_E)=(5,5)$ for C1, whose false-nearest-neighbor fraction at $d_E=5$ is $0.9\%$, $(6,5)$ for C3, and $(9,7)$ for C2, with $\tau$ in samples (Table~S5 and Fig.~S1 of the supplementary material); the value $d_E=5$ for the chaotic tasks is consistent with the sufficient condition $d_E>2D_0\simeq2D_{\mathrm{KY}}=4.55$ of Sec.~\ref{sec:attractor}.

Seven imputers are compared. Two baselines have restricted inputs by design. A static GPR uses only the current value $x_n$. GPR on the delay window of the complete channel alone is the FPRM of Ref.~\onlinecite{Wu2026FPRM}, reimplemented here with a symmetric window and, because the cost of exact GPR grows as the cube of the training-set size, a random subset of 1200 labeled samples; it is referred to as the original FPRM throughout. The remaining imputers receive the full conditioning vector of Eq.~\eqref{eq:context}: a GPR, labeled GPR (full context) in tables and figures, the cVAE, the DDPM, the neural spline flow (NSF), and the flow trained with the objective of Eq.~\eqref{eq:loss}, referred to as NSF\,+\,physics or as the physics-regularized flow. Architectures and training settings are listed in Appendix~\ref{app:hyper}.

Performance is evaluated on 1500 randomly chosen missing entries per run against the noise-free truth, with 64 samples per entry. Following Ref.~\onlinecite{Wu2026FPRM}, the posterior median $y_{\mathrm{med}}$ is scored by the Pearson correlation $\rho$ with the truth and by the normalized root-mean-square error, $\mathrm{NRMSE}=\langle(y_{\mathrm{med}}-y)^{2}\rangle^{1/2}/\mathrm{std}(y)$. Because a point estimate cannot represent an ambiguous target, the full predictive distribution $Q$ is scored by the continuous ranked probability score (CRPS)~\cite{Gneiting2007},
\begin{equation}
\mathrm{CRPS}(Q,y) = \mathbb E_Q|Y-y| - \tfrac12\,\mathbb E_Q|Y-Y'| ,
\label{eq:crps}
\end{equation}
where $Y$ and $Y'$ are independent draws from $Q$ and $y$ is the true value. The first term measures the distance of the forecast from the truth, and the second corrects for the intrinsic spread of the forecast. The score is strictly proper, i.e., its expected value is minimized only by the true conditional distribution, and it reduces to the absolute error for a point forecast; it is normalized by the standard deviation of the target. Calibration is assessed by the empirical coverage of the central $90\%$ interval~\cite{GneitingBalabdaoui2007}, and the branch ambiguity by the fraction of samples with the correct sign of the first target. In C1, the fidelity of the invariant measure is quantified by the sliced Wasserstein-1 distance $W_1$ between the true and imputed points of the stroboscopic section in $(s,c,\dot\theta)$, that is, the one-dimensional Wasserstein distance between projected point sets averaged over random projection directions~\cite{Bonneel2015}, with one posterior draw per missing section point.
